\ifx\section\undefined
  \PackageError{caption}{Use hypcap_false.tex or hypcap_true.tex}
  \endinput
\fi

\errorcontextlines=3
\sloppy

\usepackage{float}
\usepackage{floatflt}
\usepackage[rightFloats]{fltpage}
\usepackage{listings}
\usepackage{longtable}
\usepackage{picins}
\usepackage{picinpar} % must be loaded *after* picins! (why?)
\usepackage{sidecap}
\usepackage{supertabular}
\usepackage{threeparttable}
\usepackage{wrapfig}

\usepackage{hyperref}
\csname AfterHyperref\endcsname

\captionsetup{figurewithin=section,tablewithin=section}

\newcommand\sample{Some text for our page
 that is reused over and over again. }

\ifx\floatstyle\undefined \else
\floatstyle{ruled}
\newfloat{Program}{tbp}{lop}[section]
%\floatname{Program}{Programm} % german translation
\floatstyle{boxed}
\newfloat{Example}{t}{lox}[section]
%\floatname{Example}{Beispiel} % german translation
\fi

\usepackage{shortvrb}
\DeclareRobustCommand\cs[1]{\texttt{\char`\\#1}}

\newcommand{\CAPi}{%
  The dazed brown fox quickly gave 12345--67890 jumps!
  The dazed brown fox quickly gave 12345--67890 jumps!
}
  
\newcommand{\CAPii}{%
  The dazed brown fox quickly gave 12345--67890 jumps!
  The dazed brown fox quickly gave 12345--67890 jumps!
  The dazed brown fox quickly gave 12345--67890 jumps!
  The dazed brown fox quickly gave 12345--67890 jumps!
}
  
\newcommand{\SHORTCAPi}{%
  The dazed brown fox etc.}  

\begin{document}

\section{Lists}

\listoffigures
\listoftables
\ifx\listof\undefined \else
\listof{Program}{List of Programs}
\listof{Example}{List of Examples}
\fi
\ifx\lstlistoflistings\undefined \else
\lstlistoflistings
\fi

\clearpage
\section{References}
\makeatletter
\newcommand*\Ref{\@ifstar{\@Ref{\texttt{***} \cs{nameref} \textit{does not work for} \texttt}}{\@Ref\nameref}}
\newcommand*\@Ref[2]{\autoref{#2}: #1{#2} (on page \pageref{#2})}
\makeatother

\noindent
\Ref{floatingfigure.1}\\
\Ref{floatingfigure.2}\\
\Ref{FPfigure.1}\\
\Ref{figwindow.1}\\
% \nameref{} does not work with parpic & SCfigure, at least without caption package
\ifx\captionlistentry\undefined
\Ref*{parpic.1}\\
\Ref*{parpic.2}\\
\Ref*{parpic.3}\\
\Ref*{SCfigure.1}\\
\else
\Ref{parpic.1}\\
\Ref{parpic.2}\\
\Ref{parpic.3}\\
\Ref{SCfigure.1}\\
\fi
%\Ref{measuredfigure.1}\\ % hyperref anchor is ok, but reference is broken
%\Ref{measuredfigure.2}\\ %  due to an error in the threeparttable package
\Ref{wrapfigure.1}\\
 \\
\Ref{longtable.1}\\
\Ref{longtable.2}\\
\Ref{longtable.3}\\
\Ref{longtable.4}\\
\Ref{longtable.5}\\
\Ref{longtable.6}\\
\Ref{longtable.7}\\
\Ref{longtable.8}\\
% \nameref{} does not work with parpic & SCfigure, at least without caption package
\ifx\captionlistentry\undefined
\Ref*{SCtable.1}\\
\else
\Ref{SCtable.1}\\
\fi
\Ref{supertabular.1}\\
\Ref{supertabular.2}\\
\Ref{threeparttable.1}\\
\Ref{threeparttable.2}\\
 \\
\def\Programname{Program}%
\Ref{program.1}\\
\def\Examplename{Example}%
\Ref{example.1}\\
\Ref{example.2}\\
 \\
\Ref{listing.1}\\
\Ref{listing.2}\\
\Ref{list:ZeitMarkert}

\ifx\floatstyle\undefined \else
\clearpage
\section{The float package}

\begin{Program}[H]
\begin{verbatim}
#include <stdio.h>

int main(int argc, char **argv)
{
       int i;
       for (i = 0; i < argc; ++i)
               printf("argv[%d] = %s\n", i, argv[i]);
       return 0;
}
\end{verbatim}
\caption[The first program]{%
  The first program. This hasn't got anything to do with the package
  but is included as an example. Note the \texttt{ruled} float style.%
  \label{program.1}}
\end{Program}

\MakeShortVerb{\|}

\begin{Example}[H]
\begin{verse}
|\floatstyle{ruled}|\\
|\newfloat{Program}{tbp}{lop}[section]|\\
\dots\ loads o' stuff \dots\\
|\begin{Program}|\\
|\begin{verbatim}|\\
\dots\ program text \dots\\
|\end{verbatim}|\\
|\caption{|\dots\ caption \dots|}|\\
|\end{Program}|
\end{verse}
\caption[This is another silly floating Example]{%
  This is another silly floating Example. Except that this one
  doesn't actually float because it uses the {\tt[H]} optional parameter
  to appear \textbf{Here}. (Gotcha.)}\label{example.1}
\end{Example}

% Same, but with \label inside \caption:
\begin{Example}[H]
\begin{verse}
|\floatstyle{ruled}|\\
|\newfloat{Program}{tbp}{lop}[section]|\\
\dots\ loads o' stuff \dots\\
|\begin{Program}|\\
|\begin{verbatim}|\\
\dots\ program text \dots\\
|\end{verbatim}|\\
|\caption{|\dots\ caption \dots|}|\\
|\end{Program}|
\end{verse}
\caption[This is another silly floating Example]{%
  This is another silly floating Example. Except that this one
  doesn't actually float because it uses the {\tt[H]} optional parameter
  to appear \textbf{Here}. (Gotcha.)\label{example.2}}
\end{Example}

\DeleteShortVerb{\|}
\fi

\ifx\floatingfigure\undefined \else
\clearpage
\section{The floatflt package}

\begin{center}
Internet Wiretap Edition of\\[2mm]
{\large THE TELL-TALE HEART}\\
by\\ {\large EDGAR ALLAN POE}\\[2mm]
From \textit{The Works of Edgar Allan Poe: Tales Vol I}\\
J. B. Lippincott Co, Copyright 1895.\\[2mm]
This text is placed into the Public Domain (May 1993).\\
\end{center}

\noindent\textit{\Large The Tell-Tale Heart}

\noindent TRUE! nervous, very, very dreadfully nervous
I had been and am; but why WILL you say
that I am mad? The disease had sharpened
my senses, not destroyed, not dulled them. Above
all was the sense of hearing acute. I heard all things
in the heaven and in the earth. I heard many things
in hell. How then am I mad? Hearken! and observe how healthily, how
calmly, I can tell you the
whole story.

\begin{floatingfigure}{60mm}
\caption[The \texttt{floatingfigure} environment with \texttt{60mm}]%
        {The \texttt{floatingfigure} environment with \texttt{60mm}
         for \textit{width} and no \textit{option}.}
\label{floatingfigure.1}
\centering %\begin{center}
The first figure to\\ use the environment\\ \texttt{floatingfigure}
%\end{center}
\end{floatingfigure}
It is impossible to say how first the idea entered
my brain, but, once conceived, it haunted me day
and\linebreak night. Object there was none. Passion there
was none. I loved the old man. He had never
wronged me. He had never given me insult. For
his gold I had no desire. I think it was his eye!
Yes, it was this! One of his eyes resembled that
of a vulture -- a pale blue eye with a film over it.
Whenever it fell upon me my blood ran cold, and
so by degrees, very gradually, I made up my mind
to take the life of the old man, and thus rid myself
of the eye for ever.

\begin{floatingfigure}[r]{40mm}
\begin{center}
Another figure to\\ use the environment\\ \texttt{floatingfigure}
\end{center}
\caption[The \texttt{floatingfigure} environment with \texttt{40mm}]%
        {The \texttt{floatingfigure} environment with \texttt{40mm}
for \textit{width} and the \texttt{r} \textit{option}.}\label{floatingfigure.2}
\end{floatingfigure}
Now this is the point. You fancy me mad. Madmen
know nothing. But you should have seen me.
You should have seen how wisely I proceeded --
with what caution -- with what foresight, with what
dissimulation, I went to work! I was never kinder
to the old man than during the whole week before
I killed him. And every night about midnight I
turned the latch of his door and opened it oh, so
gently! And then, when I had made an opening
sufficient for my head, I put in a dark lantern all
closed, closed so that no light shone out, and then
I thrust in my head. Oh, you would have laughed
to see how cunningly I thrust it in! I moved it
slowly, very, very slowly, so that I might not
disturb the old man's sleep. It took me an hour to
place my whole head within the opening so far that
I could see him as he lay upon his bed. Ha! would
a madman have been so wise as this? And then
when my head was well in the room I undid the
lantern cautiously -- oh, so cautiously -- cautiously
(for the hinges creaked), I undid it just so much
that a single thin ray fell upon the vulture eye.
And this I did for seven long nights, every night
just at midnight, but I found the eye always closed,
and so it was impossible to do the work, for it was
not the old man who vexed me but his Evil Eye\ldots
\fi

\ifx\FPfigure\undefined \else
\clearpage
\section{The fltpage package}
\captionsetup[FPfigure]{labelfont=bf}

The following example will produce a caption below and an almost empty facing
page:
\begin{FPfigure}
 \captionsetup{textfont=it}
 \caption[A caption alone on a page]%
         {A caption alone on a page will belong to the
          following figure without a caption!}
 \label{FPfigure.1}
 \newlength\tempdima
 \setlength\tempdima\textheight
 \addtolength\tempdima{-7pt}%
 \fbox{\rule[-10cm]{0pt}{\tempdima}Frame symbolizing figure on a whole page.\hspace{5cm}}
\end{FPfigure}
\begin{quote}
  \verb|\usepackage[rightFloats]{fltpage}|\\
  \ldots\\
  \verb|\begin{FPfigure}|\\
  \verb|  \caption{A caption alone ... figure without a caption!}|\\
  \verb|  \fbox{\rule[-10cm]{0pt}{\textheight}Frame ... page.\hspace{5cm}}|\\
  \verb|\end{FPfigure}|
\end{quote}

\par\bigskip\noindent
(The caption label font should be `bf', the text font `it')

\fi

\ifx\lstlisting\undefined \else
\clearpage
\section{The listings package}

   \begin{lstlisting}[caption={Useless code},label=listing.1]
   for i:=maxint to 0 do
   begin
       { do nothing }
   end;
   \end{lstlisting}
   \begin{lstlisting}[caption={Useless code},label=listing.2,captionpos=b]
   for i:=maxint to 0 do
   begin
       { do nothing }
   end;
   \end{lstlisting}

   \begin{lstlisting}[float=htb, caption={Messung der Ausfuehrungszeit}, label=list:ZeitMarkert]
     \\Code
   \end{lstlisting}

\fi

\ifx\longtable\undefined \else
\clearpage
\section{The longtable package}

\noindent\rule{\textwidth}{1pt}% show textwidth

\begin{longtable}[l]{l}
\caption{Left aligned longtable}\label{longtable.1}\\
This is only a test\\
\end{longtable}

\begin{longtable}[l]{l}
\caption[Left aligned longtable (2)]%
        {Left aligned longtable left aligned longtable
         left aligned longtable left aligned longtable}\label{longtable.2}\\
This is only a test\\
\end{longtable}

\noindent\rule{\textwidth}{1pt}% show textwidth

\begin{longtable}[r]{l}
\caption{Right aligned longtable}\label{longtable.3}\\
This is only a test\\
\end{longtable}

\begin{longtable}[r]{l}
\caption[Right aligned longtable (2)]%
        {Right aligned longtable right aligned longtable
         right aligned longtable right aligned longtable}\label{longtable.4}\\
This is only a test\\
\end{longtable}

\noindent\rule{\textwidth}{1pt}% show textwidth

\begin{longtable}[c]{l}
\caption{Centered longtable}\label{longtable.5}\\
This is only a test\\
\end{longtable}

\begin{longtable}[c]{l}
\caption[Centered longtable (2)]%
        {Centered longtable centered longtable centered longtable
         centered longtable centered longtable}\label{longtable.6}\\
This is only a test\\
\end{longtable}

\noindent\rule{\textwidth}{1pt}% show textwidth

\begin{longtable}{c}
This is only a test\\
\caption{Longtable with `below' caption}\label{longtable.7}\\
\end{longtable}

\begin{longtable}{c}
This is only a test\\
\caption[Longtable with `below' caption (2)]%
        {Longtable with `below' caption Longtable with `below' caption
         Longtable with `below' caption Longtable with `below' caption}\label{longtable.8}\\
\end{longtable}

\noindent\rule{\textwidth}{1pt}% show textwidth
\fi

\ifx\window\undefined \else
\clearpage
\section{The picinpar package}

\font\yn=cmss17 scaled \magstep5 % or something else which is big (yinit?)
\begin{verbatim}
\begin{window}[0,l,{\yn V},{}]
or einigen Jahren wurde von ...
... Abschnitts erfolgten automatisch.
\end{window}
\end{verbatim}

\begin{window}[0,l,{\yn V},{}]
or einigen Jahren wurde von Donald E.~Knuth im TUGboat ein kleines
Problem mit der Bitte um L"osung vorgestellt. Es handelte sich darum,
in einem Paragraphen ein Fenster zu erzeugen, in das man beliebigen Text
oder eine Zeichnung hineinsetzen kann. Prompt kamen dann in den folgenden
Ausgaben L"osungsvorschl"age: Einer von DEK pers"onlich, der andere von
Alan Hoenig. Der letztgenannte brachte die elegantere L"osung, die keine
manuellen Korrekturen mehr notwendig machte. Sein Makro verlangte lediglich
in den Parametern Informationen "uber die Breite und H"ohe der
freizulassenden Stelle im Paragraphen. Die Einz"uge und der Satz der
Fragmente des Abschnitts erfolgten automatisch.
\end{window}

\begin{verbatim}
\begin{figwindow}[2,r,{
\unitlength1cm
\begin{picture}(3,1.4)
\put(0.7,0.7){\circle*{0.2}}     \put(0.7,0.7){\circle{1.2}}
\put(0.7,0.7){\vector(0,1){0.6}} \put(2.5,0.7){\circle*{0.5}}
\end{picture}
},{Kreise und Pfeile}]
Was leisten nun diese Macros ...
... sieht hierbei wie folgt aus:
\end{figwindow}
\end{verbatim}

\begin{figwindow}[2,r,{
\unitlength1cm
\begin{picture}(3,1.4)
\put(0.7,0.7){\circle*{0.2}}     \put(0.7,0.7){\circle{1.2}}
\put(0.7,0.7){\vector(0,1){0.6}} \put(2.5,0.7){\circle*{0.5}}
\end{picture}
},{Kreise und Pfeile\label{figwindow.1}}]
Was leisten nun diese Macros? Einen kleinen Eindruck hat man schon
zu Anfang dieses Artikels bekommen: der erste Buchstabe des Absatzes
ist in einer anderen Schriftgr"o"se gedruckt und in den Absatz eingepa"st.
Es mu"s aber nicht unbedingt Text sein, der so einger"uckt wird. Es kann
ebenso eine \verb?picture?--Umgebung sein, die so in den Absatz eingef"ugt
wird. Das rechts stehende Beispiel, das aus "'\LaTeX , Eine Einf"uhrung"'
bekannt sein d"urfte, ist in einer \verb?minipage?--Umgebung gesetzt
worden. Die Eingabe im Text sieht hierbei wie oben aus.
\end{figwindow}
\fi

\ifx\parpic\undefined \else
\clearpage
\section{The picins package}

\newcommand\FOR{\(\displaystyle E=mc^2\)}

\piccaptioninside
%\pichskip{1.5em}
\piccaption{Einstein's formula.\label{parpic.1}}
\parpic(50mm,10mm)[s]{\FOR}
\sample\sample\sample

\piccaptionside
\piccaption{Einstein's formula.\label{parpic.2}}
\parpic(30mm,10mm)[s]{\FOR}
\sample

\piccaptiontopside
\piccaption{Einstein's formula.\label{parpic.3}}
\parpic(30mm,10mm)[sr]{\FOR}
\sample\sample
\fi

\ifx\SCfigure\undefined \else
\clearpage
\section{The sidecap package}

\newcommand{\FIG}[2][]{%
  \begingroup
    \def\xxx{#1}%
    \ifx\xxx\empty
      \setlength{\unitlength}{\linewidth}%
      \addtolength{\unitlength}{-2\fboxrule}%
      \setlength{\unitlength}{.1\unitlength}%
    \else  
      \setlength{\unitlength}{#1}%
    \fi% 
    \setlength{\fboxsep}{0pt}%
    \ifcase#2\relax
    \or%
      \fbox{%
	\begin{picture}(4,6)%
	  \put(1,5){\circle{1}}%
	  \put(3,5){\circle{1}}%
	  \put(2,3){\circle{1}}%
	  \put(1,1){\circle{1}}%
	  \put(3,1){\circle{1}}%     
	\end{picture}}%
    \or%	
      \fbox{%
	\begin{picture}(10,4)%
	  \put(1,3){\circle{1}}%
	  \put(9,3){\circle{1}}%
	  \put(5,2){\circle{1}}%
	  \put(1,1){\circle{1}}%
	  \put(9,1){\circle{1}}%     
	\end{picture}}%
    \or% (doesn't make a difference to use [b] or [t] or [c] here !!!)
      \begin{minipage}{5\unitlength}%
        text text text text text text text text
	text text text text text text text text
	text text text text text text text text
	text text text text text text text text
	text text text text text text text text
	text text text text text text text text
	text text text text text text text text
	text text text text text text text text.
      \end{minipage}%	
    \or% (doesn't make a difference to use [b] or [t] or [c] here !!!)
      \begin{minipage}{5\unitlength}%
        \strut
        text text text text text text text text
	text text text text text text text text
	text text text text text text text text
	text text text text text text text text
	text text text text text text text text
	text text text text text text text text
	text text text text text text text text
	text text text text text text text text.
	\unskip\strut
      \end{minipage}%	
    \fi	
  \endgroup  
}

\begin{SCfigure}[][!hb]
  \FIG[1cm]{1}%
  \caption[\SHORTCAPi]{\CAPi}
  \label{SCfigure.1}
\end{SCfigure}

\begin{SCtable}[][!hb]
  \caption[\SHORTCAPi]{\CAPi}  
  \FIG[1cm]{1}%
  \label{SCtable.1}
\end{SCtable}

\begin{SCfigure}[][!hb]
  \FIG[0.5cm]{1}%
  \caption*{\SHORTCAPi}
% \label{SCfigure.2}
\end{SCfigure}
\fi

\ifx\supertabular\undefined \else
\clearpage
\section{The supertabular package}

\begin{center}
\tablecaption{The ISOGRK3 entity set (tablecaption)\label{supertabular.1}}
\begin{supertabular}{lll}
\bfseries Entity&\bfseries  Unicode Name&\bfseries  Unicode\\
\hline
alpha              & GREEK SMALL LETTER ALPHA            & 03B1\\
beta               & GREEK SMALL LETTER BETA             & 03B2\\
chi                & GREEK SMALL LETTER CHI              & 03C7\\
\end{supertabular}
\end{center}

\begin{center}
\bottomcaption{The ISOGRK3 entity set (bottomcaption)}\label{supertabular.2}
\begin{supertabular}{lll}
\bfseries Entity&\bfseries  Unicode Name&\bfseries  Unicode\\
\hline
alpha              & GREEK SMALL LETTER ALPHA            & 03B1\\
beta               & GREEK SMALL LETTER BETA             & 03B2\\
chi                & GREEK SMALL LETTER CHI              & 03C7\\
\end{supertabular}
\end{center}
\fi

\ifx\threeparttable\undefined \else
\clearpage
\section{The threeparttable package}

\hrule

\begin{center}
\begin{threeparttable}
\caption{The Skewing Angles ($\beta$) for $\fam0 Mu(H)+X_2$ and
   $\fam0 Mu(H)+HX$~\tnote{a}}
   \label{threeparttable.1}
\begin{tabular}{rlcc}
\hline
&   & $\fam0 H(Mu)+F_2$ & $\fam0 H(Mu)+Cl_2$ \\
\hline
&$\beta$(H)  & $80.9^\circ\tnote{b}$ & $83.2^\circ$ \\
&$\beta$(Mu) & $86.7^\circ$ & $87.7^\circ$ \\
\hline
\end{tabular}
\begin{tablenotes}
\item[a] for the abstraction reaction, $\fam0 Mu+HX \rightarrow MuH+X$.
\item[b] 1 degree${} = \pi/180$ radians.
\end{tablenotes}
\end{threeparttable}
\end{center}

\hrule

\begin{center}
\begin{threeparttable}
\begin{tabular}{rlcc}
\hline
&   & $\fam0 H(Mu)+F_2$ & $\fam0 H(Mu)+Cl_2$ \\
\hline
&$\beta$(H)  & $80.9^\circ\tnote{b}$ & $83.2^\circ$ \\
&$\beta$(Mu) & $86.7^\circ$ & $87.7^\circ$ \\
\hline
\end{tabular}
\begin{tablenotes}
\item[a] for the abstraction reaction, $\fam0 Mu+HX \rightarrow MuH+X$.
\item[b] 1 degree${} = \pi/180$ radians.
\end{tablenotes}
\caption{The Skewing Angles ($\beta$) for $\fam0 Mu(H)+X_2$ and
   $\fam0 Mu(H)+HX$~\tnote{a}}
   \label{threeparttable.2}
\end{threeparttable}
\end{center}

\hrule

\begin{center}
\begin{measuredfigure}
\begin{tabular}{rlcc}
\hline
&   & $\fam0 H(Mu)+F_2$ & $\fam0 H(Mu)+Cl_2$ \\
\hline
&$\beta$(H)  & $80.9^\circ\tnote{b}$ & $83.2^\circ$ \\
&$\beta$(Mu) & $86.7^\circ$ & $87.7^\circ$ \\
\hline
\end{tabular}
\caption{The Skewing Angles ($\beta$) for $\fam0 Mu(H)+X_2$ and
   $\fam0 Mu(H)+HX$~\tnote{a}}
%%%   \label{measuredfigure.1}
\end{measuredfigure}
\end{center}

\hrule

\begin{center}
\begin{measuredfigure}
\caption{The Skewing Angles ($\beta$) for $\fam0 Mu(H)+X_2$ and
   $\fam0 Mu(H)+HX$~\tnote{a}}
%%%   \label{measuredfigure.2}
\begin{tabular}{rlcc}
\hline
&   & $\fam0 H(Mu)+F_2$ & $\fam0 H(Mu)+Cl_2$ \\
\hline
&$\beta$(H)  & $80.9^\circ\tnote{b}$ & $83.2^\circ$ \\
&$\beta$(Mu) & $86.7^\circ$ & $87.7^\circ$ \\
\hline
\end{tabular}
\end{measuredfigure}
\end{center}

\hrule
\fi

\ifx\wrapfigure\undefined \else
\clearpage
\section{The wrapfig package}

The starting place for the wrapfigure
environment was manually determined in
the current ex-
\begin{wrapfigure}[7]{l}[0.2\width]{0pt}
  \centering
  \fbox{This is a ``wrapfigure''.}
  \caption{An example of the
    \texttt{wrapfigure} environment}
  \label{wrapfigure.1}
\end{wrapfigure}
sample by first setting the text without
the figure to find the linebreaks.

\sample \sample \sample
\fi

\clearpage
That's all folks!
\end{document}
